\section{Introduction}

Neural networks that apply one block again and again, in a loop, have learned to solve Sudoku, mazes and the textbook
algorithms of the CLRS-30 benchmark\cite{wang2025hrm,jolicoeur2025trm,velickovic2022clrs,bansal2022end}. They are
trained through their loops, and that is slow. On Sudoku-Extreme the strongest of them train for \nm{looped.fit.lo} to
\nm{looped.fit.hi} hours on a GPU\cite{hoover2026classical}. The looped attention blocks of our earlier paper needed
about an hour of GPU time for each CLRS-30 task\cite{mohsenvand2026zero}.

Hoover, Blayney and Whittington showed that much of this cost can be avoided\cite{hoover2026classical}. Their model,
CRAM, is the pairwise energy that Hopfield and Tank proposed for puzzles\cite{hopfield1985neural}, with its tables fitted
by pseudolikelihood\cite{besag1975statistical}: each cell of a solved board is predicted from the others, so training
needs no loop and no search. CRAM trains in \nm{cram.sudoku.fit} milliseconds and solves \nm{cram.sudoku.pct}\% of
Sudoku-Extreme. Its authors name its two limits. Its costs belong to particular cells, so a model fitted at one board
size says nothing at another. And it has no place for a quantity the answer does not show: on Maze-Hard it needs a
hidden distance layer designed for the task, and learning such quantities in general is left open.

What a fit at a solved example learns is a check, a rule that each part of the answer satisfies given the rest. A check
is not yet an algorithm. At a finished answer every fact is in place, so facts may hold one another up in a circle, and
nothing in the fit asks whether the rule, started from nothing, would put them there. Logic programming names the gap
exactly: a finished answer the rules accept is a supported model; what the rules derive from nothing is their least
model; and the two coincide when the support has a direction\cite{clark1978negation,fages1994consistency}.

We give the check that direction, and it becomes the algorithm. We start from the cold layer of our earlier
paper\cite{mohsenvand2026zero} (Fig.~\ref{fig:arch}). Its state lives on pairs of items. It judges each pair $(a,b)$ through every third item
$k$, by the legs $(a,k)$ and $(k,b)$, and reads these witnesses at zero temperature: whether one exists, and how two
counts compare. Its rules attach to relations, never to positions, so they have no size in them. We fit its check at
finished answers, CRAM's way, in three steps.
\begin{itemize}
\setlength{\itemsep}{0pt}
\item \emph{Complete the answer from itself.} Closures and counts of the answer's own relations add what the answer
hides; a parent pointer, for instance, gives each node its depth.
\item \emph{Fit the check and keep every condition that held on all training answers.} Among them is the level
condition, that every part rests on parts placed before it. A finished answer never contradicts it, so a fit alone can
drop it.
\item \emph{Run the check cold.} From an empty state, apply the check and recompute the completion until nothing changes.
Where nothing is forced, as on hard puzzles, a search uses the check as its stop.
\end{itemize}

The check then trains like CRAM and runs like our looped layer (Fig.~\ref{fig:idea}). All thirty CLRS-30 algorithms are
exact on the official test and on every held-out input, up to \nm{held.max.nodes} nodes and \nm{held.max.items} items.
Fifteen of them, among them breadth-first search, the minimum spanning trees of Kruskal and Prim, longest common
subsequence and both string matchers, are learned from their final answers with comparisons and counts alone, each
fitted in \nm{clrs.g1.fit.lo} to \nm{clrs.g1.fit.hi} on one laptop core; the other fifteen read one quantity we supply,
such as the sum of a distance and an edge weight. The fitted rules read like the textbook: a node's breadth-first parent
is its lowest-numbered neighbour one level up. On CRAM's benchmarks, one rule fitted in about 20 milliseconds keeps
exactly the row, column and box relations of Sudoku. It solves all \nm{sudoku.ntest} Sudoku-Extreme puzzles with no
wrong answer, with \nm{sudoku.ops.ratio} times fewer operations than CRAM, and solves larger boards unchanged. A maze
rule that reads the cold run's own clock, in place of a distance layer built for the task, solves every Maze-Hard maze,
and every maze we built on which CRAM's released model leaves the shortest path. The fitted rules are small tables with
no size in them, so finite checks prove the puzzle rules, and the checks of several algorithms, correct at every size
(Fig.~\ref{fig:cost} sets cost against reach). Read the same way, CRAM's released weights show its hidden distance
layer cooling, as it trains, towards the hard minimum that our cold layer reads.
